% Methods section for the TZ Fornacis joint fit. Numbers are those of the
% September 2026 analysis (tzfor_teff2_revision.ipynb, sections 7.3-7.5).

\section{Joint fitting of TZ Fornacis}
\label{sec:tzfor_fit}

\subsection{Data}

We fit two data sets simultaneously. The photometry is the Str\"omgren $b$ and
$y$ light curves of \citet{Clausen1991}, 1169 epochs in each band (2338 points),
with per-point uncertainties of 4.1 and 3.5\,mmag taken from the out-of-eclipse
scatter. The spectroscopy is 21 HARPS exposures covering 4450--6750\,\AA, which we
continuum-normalise with SuppNet and resample onto a common logarithmic
wavelength grid of 156\,133 pixels per epoch, $3.0\times10^{6}$ pixels in all.
Pixels affected by cosmic rays, order edges and the H$\alpha$ core, whose LTE
profile is known to be inadequate, are excluded from the outset. The two stars'
line systems are separated by up to 80\,km\,s$^{-1}$ over the orbit, and the
exposures sample the full radial-velocity cycle, though none falls inside an
eclipse.

\subsection{Forward model}

\paragraph{Geometry.}
The stellar surfaces are Roche meshes generated with PHOEBE \citep{Prsa2016} at
the literature geometry: the masses, radii and inclination of
\citet{Andersen1991} and the period of \citet{Gallenne2016}. Meshes of 1718 and
1712 elements are built at each of the 21 spectroscopic epochs, and coarser meshes
of 1372 and 1414 elements at 56 orbital phases concentrated on the two eclipses
for the photometry. Each element carries its projected area, its occluded
fraction, its direction cosine $\mu$, its surface gravity and its line-of-sight
velocity, which combines the orbital motion with rigid rotation at
5.5\,km\,s$^{-1}$ for the giant (synchronous) and 38\,km\,s$^{-1}$ for the
secondary \citep{Andersen1991}.

\paragraph{Atmospheres.}
Specific intensities come from a Transformer-Payne emulator trained on MARCS
model atmospheres over the HARPS range, which returns $I_\lambda(\mu)$ and the
corresponding continuum for a given $T_{\rm eff}$, $\log g$, [Fe/H],
microturbulence and the abundances of $\alpha$ elements, C, N and O. Gravity
darkening is applied element by element as
$T_{\rm eff,i}=T_{\rm eff}\,10^{\beta(\log g_i-\overline{\log g})}$ with
$\beta=0.079$; the exponent is not constrained by TZ~For, whose intra-star
temperature spread is 3\,K and 16\,K for the two components, and is held fixed.
Microturbulence is fixed at 1.5\,km\,s$^{-1}$ and the abundance ratios at
solar-scaled values; freeing $\alpha$, C, N, O for the giant moves them by
$\le0.05$\,dex and changes the temperatures by less than the systematic error
quoted below.

\paragraph{Synthesis.}
Evaluating the emulator for every element at every wavelength is exact but costs
2.5\,min per gradient on an A100 for a 12-window subset alone, which rules it out
inside an optimiser. We instead use a synthesis kernel: the elements are binned
into 16 nodes in $\mu$ and two in $\log g$, and each element's projected area is
accumulated, at its Doppler shift, into a velocity kernel on a wavelength grid
oversampled twice relative to the data. The emulator is evaluated once per node
and convolved with the kernel by FFT. This is a discretisation of the
per-element sum rather than a different physical model, and the two agree to a
maximum of $8\times10^{-4}$ and an rms of $1\times10^{-4}$ in normalised flux
across the full range, two orders of magnitude below the residuals of the fit.
The joint gradient then costs 36\,s. The final model spectra and light curves
shown in the figures are recomputed with the exact per-element synthesis at the
fitted parameters.

\paragraph{Broadening and blending.}
Each star's line and continuum channels are broadened by a Gaussian combining the
instrumental resolution ($R=115\,000$) with a free macroturbulent velocity,
$v_{\rm mac,1}$ and $v_{\rm mac,2}$. The projected rotation of each star is
scaled by a free factor $s_{1,2}$ applied to the element velocities about their
orbital mean, so that the kernels are rebuilt inside the loss. The blended
spectrum is the ratio of summed line to summed continuum channels,
$(F_1+F_2)/(C_1+C_2)$, so that the line-depth dilution of each component follows
from the light ratio rather than being a free parameter. A systemic-velocity
offset of 0.25\,km\,s$^{-1}$ relative to the mesh velocities is applied as a
shift of the model wavelength grid.

\paragraph{Continuum.}
The residual normalisation of each observed spectrum is absorbed by a cubic
B-spline with 40\,\AA\ knot spacing, multiplied into the model. Its coefficients
are the weighted least-squares solution at every evaluation of the loss, so they
are profiled out rather than sampled, and the fitted parameters cannot trade
against them.

\paragraph{Photometry.}
Passband magnitudes are integrated from the model flux on a 16\,000-point grid
between 4300 and 5900\,\AA; coarser sampling of the line-rich emulator spectra
biases the integrals by several mmag. The model light curve is interpolated in
phase to the observation times with a free phase offset $\Delta\varphi$ that
absorbs the difference between the photometric and the mesh ephemerides, and a
zero point per band is set by the median residual.

\subsection{Objective and optimisation}

The objective is the sum of the photometric and spectroscopic $\chi^2$, the
former with the per-band uncertainties above, the latter with a per-epoch,
per-block noise estimate from the HARPS pipeline added in quadrature to a 0.5\%
floor. The free parameters are $T_{\rm eff,1}$, $T_{\rm eff,2}$, [Fe/H],
$\Delta\varphi$, $v_{\rm mac,1}$, $v_{\rm mac,2}$, $s_1$ and $s_2$. The forward
model is written in JAX, so the objective and its gradient with respect to all
parameters are obtained from a single reverse-mode pass; we minimise with
L-BFGS-B on scaled variables (100\,K, 0.1\,dex, 0.001 in phase, 1\,km\,s$^{-1}$
and 0.1 in the rotation scale) within physical bounds. Two numerical details
matter at this size. The emulator's matrix products must run at full
single-precision rather than the tensor-core default, which otherwise makes the
objective rough at the 10\,K scale and stalls line searches; and a fit that
stalls in its first line search is restarted from a perturbed point. Fits
converge in 20--30 iterations.

\subsection{The emulator's line-list floor}

Stacking the residuals of all epochs in each star's rest frame shows that the
2--3\% rms that remains after the fit is a fixed line-by-line pattern in the
giant's frame: some 960 features coherent across the 21 epochs, 285 of them above
10\%, roughly equally split between lines the emulator makes too deep and too
shallow. They respond to [Fe/H] and microturbulence but not to the C, N, O or
$\alpha$ abundances, and they are the line data of a 4900\,K giant rather than
its composition. We treat them in two independent ways and report both.

\paragraph{Learned mask.}
Pixels whose stacked residual exceeds 5\% at more than four times its standard
error are excluded, with two pixels of padding, in whichever star's frame they
occur. This removes 9.4\% of the pixels and lowers the spectroscopic $\chi^2$
per point from 9.5 to 3.8. To guard against circularity the mask is derived on
the odd-numbered epochs alone and applied to a fit of the even-numbered ones;
the solution is unchanged to within 2\,K and 0.003\,dex.

\paragraph{Line-strength map.}
Alternatively, a multiplicative correction $\delta(\lambda)$ to the giant's line
channel is fitted once, at fixed stellar parameters, on the odd epochs, with a
Gaussian prior of 3\% and a second-difference smoothness term. Fitting $\delta$
jointly with the stellar parameters is degenerate and is not done. Held fixed on
the even epochs, which did not inform it, the map brings the spectroscopic
$\chi^2$ per point to 1.5 and the residual to 0.8--1.2\%, and shifts the solution
by $-25$\,K, $-75$\,K and $-0.10$\,dex relative to the masked fit. That shift is
the dominant systematic in the result.

\subsection{Uncertainties}

Three statements of uncertainty are reported. A block jackknife divides the
spectrum into six 380\,\AA\ blocks and refits with each removed in turn; the
spread of the six solutions measures the sensitivity to which lines are used.
A Hamiltonian Monte Carlo posterior is drawn with NUTS on a version of the model
in which each star's disc spectra are precomputed on a $17\times7$ grid in
($T_{\rm eff}$, [Fe/H]) and interpolated with cubic splines, accurate to
$10^{-4}$, with a free noise-inflation factor per spectral block; four chains of
1500 samples after 400 warm-up steps give $\hat R\le1.001$. With $3\times10^{6}$
pixels the posterior widths are about 1\,K, and we do not quote them as the
error: they measure the statistical precision of a model whose residuals are
dominated by correlated line-list systematics. The quoted error is the jackknife
spread combined with the spread across the mask and line-strength treatments.

\subsection{Free geometry and the photometric cross-check}

To test whether the geometry rather than the atmospheres could be at fault, the
radii and inclination are also freed. The photometric model is then built inside
the loss from SPICE's own binary, two icospheres of 2000 vertices each on the
Keplerian orbit with SPICE's polygon-clipping occlusion, at 51 phases
concentrated on the eclipses, and the spectral kernels scale with the projected
areas. At the literature geometry this model reproduces the PHOEBE-mesh light
curves to 0.6\,mmag, and its occlusion matches the analytic overlap of two discs
to 0.1\%. The joint gradient rises to 48\,s.

The same fits are put to PHOEBE with its Roche meshes, ck2004 atmospheres and
irradiation. Its own optimiser, with both radii, the inclination and
$T_{\rm eff,2}$ free, returns the literature radius ratio, and its photometric
$\chi^2$ scanned across the ratio at fixed radius sum rises monotonically away
from it. Roche and rotational distortion cannot account for the difference: the
surfaces depart from spheres by 0.08\% and 0.76\% in radius, but the projected
areas, which set the light ratio, by only $+0.02$\% and $-0.13$\%. The
free-geometry solution is therefore reported as a diagnostic of the
spectroscopic light-ratio tension rather than as a radius measurement.

\subsection{Cost}

A joint fit of the full HARPS range with photometry converges in 20--30
minutes on one A100; the jackknife adds six refits of the same order, the NUTS
posterior two hours, and the exact per-element synthesis of the final model
half an hour. All of it is driven by the same differentiable forward model.
